\documentclass[10pt,a4paper]{amsart}
\usepackage[margin=19mm]{geometry}
\usepackage{amsmath,amssymb,mathtools}
\usepackage[scr=rsfs,cal=euler]{mathalfa}
\usepackage{xfrac}
\usepackage[unicode,hidelinks,bookmarks=false]{hyperref}
\newcommand{\ie}{i.e.\ }
\newcommand{\cf}{cf.\ }
\newcommand{\resp}{respectively\ }
\newcommand{\Subsection}{\subsection}
\providecommand{\nobreakdash}{\nobreak\hbox{-}}
\setlength{\emergencystretch}{2em}
\multlinegap0pt
\usepackage{amssymb}
\usepackage{float}
\usepackage{tikz}
\newtheorem{theorem}{Theorem}
\newcommand{\spec}{\operatorname{Spec}}
\title{Degree-based weighted adjacency matrices: spectra, integrality, and edge deletion effects}
\author{{\small Bilal Ahmad Rather$^{a}$ and Hilal Ahmad Ganie$^{b}$}}
\date{Source: arXiv:2603.08895v2 (CC BY 4.0)}
\begin{document}
\maketitle

\noindent\emph{Source and licence.} Degree-based weighted adjacency matrices: spectra, integrality, and edge deletion effects. Version arXiv:2603.08895v2 (CC BY 4.0), distributed by arXiv under the Creative Commons Attribution 4.0 International licence, http://creativecommons.org/licenses/by/4.0/, as recorded in the arXiv licence metadata for this exact version and shown on the abstract page https://arxiv.org/abs/2603.08895v2. Full licence notice: \texttt{licenses/arxiv-2603.08895-CC-BY-4.0.txt}.\par\smallskip

\noindent\textit{Editorial scope.} Counted unit: the article's theorem thm:spectrum\_complete\_multipartite\_updated (source0.tex lines 249--313). The article's complete original proof of it is delivered with it (source0.tex lines 315--680, one \texttt{proof} environment of 366 lines). No mathematics is written, repaired or re-derived: the statement and the proof are the article's own lines, in the article's own notation, and no retained line is substituted or re-flowed.  No further source-local context is needed: every cross-reference of every retained line resolves inside the two delivered spans.  Nothing was omitted from any retained span, so the package carries no omission record.  No retained line contains a citation, so the wrapper carries no bibliography and no bibliography entry.  No retained line contains a figure environment, an image-inclusion command or drawing code, so no image is delivered and the package declares no asset dependency.  Editorial formatting only, in addition: an AMS \texttt{amsart} preamble that restates the article's own declarations and macros used by the retained lines, namely the environments \texttt{theorem} on separate counters, together with the standard packages those lines need.  The retained environments print in this document as Theorem 1 in order of appearance, whereas the article prints the same items with the numbering of its own full document.  The complete native source of the article as distributed in the arXiv e-print for this exact version is bundled unchanged as \texttt{sources/196036/source0.tex}.
\begin{theorem}\label{thm:spectrum_complete_multipartite_updated}
	Let $G\cong K_{p_1,p_2,\ldots,p_t}$, where $t\geq 2$, $p_i\geq 1$,
	$
	n=\sum_{i=1}^t p_i,
	$
	and let
	$
	d_i=n-p_i
	$
	be the common degree of the vertices belonging to the $i$-th partite set. Let
	$
	M=(m_{ij})_{t\times t},
	$
	where
	$
	m_{ii}=0
	$
	and
	$
	m_{ij}=p_j\phi(d_i,d_j)~
	 (i\neq j).	$
	Let
	$
	D=\operatorname{diag}(p_1,\ldots,p_t)
	$
	and define
	$
	S=D^{1/2}MD^{-1/2}=(s_{ij})_{t\times t}.
	$
	Then
	$
	s_{ii}=0
	$
	and, for $i\neq j$,
	$
	s_{ij}=\sqrt{p_ip_j}\,\phi(d_i,d_j).
	$
	In particular, under the standing assumption that $\phi$ is symmetric, $S$ is a real symmetric matrix and $M$ is similar to $S$. Moreover,
	$
	\chi_{\dot A(G)}(x)
	=
	x^{n-t}\det(xI_t-S).
	$
	Consequently,
	$
	\spec(\dot A(G))
	=
	\{0^{[n-t]}\}\uplus\spec(S)
	=
	\{0^{[n-t]}\}\uplus\spec(M),
	$
	where $\uplus$ denotes multiset union. Furthermore,
	$
	\operatorname{rank}\dot A(G)
	=
	\operatorname{rank}S
	$
	and
	$
	\operatorname{mult}_{\dot A(G)}(0)
	=
	n-t+\dim\ker S.
	$
	Hence the multiplicity of $0$ is exactly $n-t$ if and only if $S$ is nonsingular.
\end{theorem}

\begin{proof}
	Let
	$
	V(G)=V_1\sqcup V_2\sqcup\cdots\sqcup V_t,
	~ |V_i|=p_i,
	$
	be the canonical multipartite decomposition of $G$. Since $G$ is complete multipartite, no two vertices belonging to the same partite set are adjacent, whereas every vertex of $V_i$ is adjacent to every vertex of $V_j$ whenever $i\neq j$. Every vertex of $V_i$ has degree
	$
	d_i=n-p_i.
	$
	Therefore all edges joining $V_i$ and $V_j$ receive the common weight
	$
	\phi(d_i,d_j).
	$
	With the vertices ordered according to $V_1,\ldots,V_t$, the matrix $\dot A(G)$ has block form
	$$
	\dot A(G)
	=
	\begin{pmatrix}
		0_{p_1}
		&
		\phi(d_1,d_2)J_{p_1\times p_2}
		&
		\cdots
		&
		\phi(d_1,d_t)J_{p_1\times p_t}
		\\
		\phi(d_2,d_1)J_{p_2\times p_1}
		&
		0_{p_2}
		&
		\cdots
		&
		\phi(d_2,d_t)J_{p_2\times p_t}
		\\
		\vdots
		&
		\vdots
		&
		\ddots
		&
		\vdots
		\\
		\phi(d_t,d_1)J_{p_t\times p_1}
		&
		\phi(d_t,d_2)J_{p_t\times p_2}
		&
		\cdots
		&
		0_{p_t}
	\end{pmatrix}.
	$$
	
	For each $i\in\{1,\ldots,t\}$, define
	$$
	U_i
	=
	\left\{
	x\in\mathbb R^{V(G)}:
	\operatorname{supp}(x)\subseteq V_i,\ 
	\sum_{u\in V_i}x_u=0
	\right\}.
	$$
	The condition that the coordinates of $x$ on $V_i$ sum to zero imposes one independent linear relation on a $p_i$-dimensional coordinate space. Hence
	$
	\dim U_i=p_i-1.
	$
	
	We first show that every vector in $U_i$ belongs to the kernel of $\dot A(G)$. Let $x\in U_i$. If $v\in V_i$, then $x$ is supported entirely on $V_i$, while there are no edges inside $V_i$. Consequently,
	$
	(\dot A(G)x)_v=0.
	$
	Now let $v\in V_j$ with $j\neq i$. Since $v$ is adjacent to every vertex of $V_i$ and all these edges have weight $\phi(d_j,d_i)$,
	$$
	(\dot A(G)x)_v
	=
	\sum_{u\in V_i}\phi(d_j,d_i)x_u
	=
	\phi(d_j,d_i)\sum_{u\in V_i}x_u
	=
	0.
	$$
	Thus every coordinate of $\dot A(G)x$ is zero, and therefore
	$
	\dot A(G)x=0.
	$
	Hence
	$
	U_i\subseteq\ker\dot A(G)
	$
	for every $i$.
	
	The subspaces $U_1,\ldots,U_t$ have pairwise disjoint supports. In particular, they are mutually orthogonal. Therefore their sum is direct, and if
	$$
	U=\bigoplus_{i=1}^t U_i,
	$$
	then
	$
	U\subseteq\ker\dot A(G)
	$
	and
	$$
	\dim U
	=
	\sum_{i=1}^t(p_i-1)
	=
	\sum_{i=1}^t p_i-t
	=
	n-t.
	$$
	This already explains the forced occurrence of at least $n-t$ zero eigenvalues.
	
	To identify the remaining spectral information, define the normalized characteristic vector of $V_i$ by
	$$
	z_i
	=
	\frac{1}{\sqrt{p_i}}\mathbf 1_{V_i},
	\qquad i=1,\ldots,t,
	$$
	and let
	$
	W=\operatorname{span}\{z_1,\ldots,z_t\}.
	$
	Since the sets $V_i$ are pairwise disjoint,
	$
	\langle z_i,z_j\rangle=0$ for $i\neq j,$
	while
	$$
	\|z_i\|^2
	=
	\frac{1}{p_i}\langle\mathbf 1_{V_i},\mathbf 1_{V_i}\rangle
	=
	1.
	$$
	Thus
	$
	z_1,\ldots,z_t
	$
	form an orthonormal basis of $W$, and
	$
	\dim W=t.
	$
	
	Next, we verify that $U$ and $W$ are orthogonal complements. If $x\in U_i$, then
	$$
	\langle x,z_i\rangle
	=
	\frac{1}{\sqrt{p_i}}
	\sum_{u\in V_i}x_u
	=
	0.
	$$
	Moreover, for $j\neq i$, the supports of $x$ and $z_j$ are disjoint, so
	$
	\langle x,z_j\rangle=0.
	$
	Hence
	$
	U\perp W.
	$
	Since
	$
	\dim U+\dim W=(n-t)+t=n,
	$
	we obtain the orthogonal decomposition
	$
	\mathbb R^{V(G)}=U\oplus W.
	$
	
	It remains to determine the action of $\dot A(G)$ on $W$. Fix $j\in\{1,\ldots,t\}$. Since $z_j$ is supported on $V_j$, there is no contribution from the diagonal block corresponding to $V_j$. For $i\neq j$, every coordinate of $\dot A(G)z_j$ belonging to $V_i$ equals
	$$
	\phi(d_i,d_j)
	\sum_{u\in V_j}\frac{1}{\sqrt{p_j}}
	=
	\phi(d_i,d_j)\frac{p_j}{\sqrt{p_j}}
	=
	\sqrt{p_j}\,\phi(d_i,d_j).
	$$
	Therefore the $V_i$-component of $\dot A(G)z_j$ is
	$
	\sqrt{p_j}\,\phi(d_i,d_j)\mathbf 1_{V_i}.
	$
	Since
	$
	\mathbf 1_{V_i}=\sqrt{p_i}\,z_i,
	$
	this can be written as
	$
	\sqrt{p_ip_j}\,\phi(d_i,d_j)z_i.
	$
	Hence
	$$
	\dot A(G)z_j
	=
	\sum_{\substack{i=1\\i\neq j}}^t
	\sqrt{p_ip_j}\,\phi(d_i,d_j)z_i.
	$$
	In particular, $\dot A(G)z_j\in W$, so $W$ is invariant under $\dot A(G)$.
	 Relative to the orthonormal basis
	$
	z_1,\ldots,z_t,
	$
	the matrix of the restriction $\dot A(G)|_W$ is 
	$
	S=(s_{ij})_{t\times t},
	$
	where
	$
	s_{ii}=0
	$
	and
	$$
	s_{ij}
	=
	\sqrt{p_ip_j}\,\phi(d_i,d_j)
	\qquad(i\neq j).
	$$
	Since $\phi$ is symmetric,
	$
	\phi(d_i,d_j)=\phi(d_j,d_i),
	$
	and consequently
	$$
	s_{ij}
	=
	\sqrt{p_ip_j}\,\phi(d_i,d_j)
	=
	\sqrt{p_jp_i}\,\phi(d_j,d_i)
	=
	s_{ji}.
	$$
	Thus $S$ is real symmetric.
	
	Choose an orthonormal basis of each $U_i$, combine these bases to obtain an orthonormal basis of $U$, and append the vectors
	$
	z_1,\ldots,z_t.
	$
	With respect to the resulting orthonormal basis of $\mathbb R^{V(G)}$, the operator $\dot A(G)$ has block diagonal matrix
	$
	0_{n-t}\oplus S.
	$
	Equivalently, there exists an orthogonal matrix $P$ such that
	$$
	P^{\top}\dot A(G)P
	=
	\begin{pmatrix}
		0_{n-t}&0\\
		0&S
	\end{pmatrix}.
	$$
	Thus $\dot A(G)$ is orthogonally similar to $0_{n-t}\oplus S$. It follows that
	$$
	\chi_{\dot A(G)}(x)
	=
	\det\left(xI_n-(0_{n-t}\oplus S)\right).
	$$
	Using the determinant of a block diagonal matrix,
	$$
	\chi_{\dot A(G)}(x)
	=
	\det(xI_{n-t})\det(xI_t-S)
	=
	x^{n-t}\det(xI_t-S).
	$$
	This proves the asserted characteristic-polynomial factorization.
	
	The same orthogonal decomposition gives the spectrum, including multiplicities:
	$$
	\spec(\dot A(G))
	=
	\{0^{[n-t]}\}\uplus\spec(S).
	$$
	Notice that the multiset notation is important here. If $S$ itself is singular, additional zero eigenvalues arise from $\spec(S)$ and must be added to the forced multiplicity $n-t$.
	
	We now connect $S$ with the equitable quotient $M$. Since
	$
	D^{1/2}
	=
	\operatorname{diag}(\sqrt{p_1},\ldots,\sqrt{p_t})
	$
	and
	$$
	D^{-1/2}
	=
	\operatorname{diag}
	\left(
	\frac{1}{\sqrt{p_1}},\ldots,
	\frac{1}{\sqrt{p_t}}
	\right),
	$$
	the $(i,j)$-entry of $D^{1/2}MD^{-1/2}$, for $i\neq j$, is
	$$
	\sqrt{p_i}
	\left(
	p_j\phi(d_i,d_j)
	\right)
	\frac{1}{\sqrt{p_j}}
	=
	\sqrt{p_ip_j}\,\phi(d_i,d_j).
	$$
	The diagonal entries remain zero. Therefore
	$
	D^{1/2}MD^{-1/2}=S.
	$
	As every $p_i$ is positive, $D^{1/2}$ is invertible. Hence $M$ and $S$ are similar, and consequently
	$
	\spec(M)=\spec(S).
	$
	This implies that
	$$
	\spec(\dot A(G))
	=
	\{0^{[n-t]}\}\uplus\spec(S)
	=
	\{0^{[n-t]}\}\uplus\spec(M).
	$$
	
	The rank formula is now immediate from the orthogonal block decomposition. Since
	$
	\dot A(G)\sim 0_{n-t}\oplus S,
	$
	similarity preserves rank and
	$
	\operatorname{rank}(0_{n-t}\oplus S)
	=
	\operatorname{rank}S.
	$
	Therefore
	$
	\operatorname{rank}\dot A(G)
	=
	\operatorname{rank}S.
	$
	
	Likewise,
	$
	\ker(0_{n-t}\oplus S)
	=
	\mathbb R^{n-t}\oplus\ker S.
	$
	Consequently,
	$
	\dim\ker\dot A(G)
	=
	n-t+\dim\ker S.
	$
	Since $\dot A(G)$ is real symmetric, it is diagonalizable, and hence the dimension of its kernel equals the algebraic multiplicity of the eigenvalue $0$. Thus
	$$
	\operatorname{mult}_{\dot A(G)}(0)
	=
	n-t+\dim\ker S.
	$$
	It follows that
	$
	\operatorname{mult}_{\dot A(G)}(0)=n-t
	$
	if and only if
	$
	\dim\ker S=0.
	$
	For a square matrix this is equivalent to
	$
	\det S\neq0,
	$
	that is, to the nonsingularity of $S$. This completes the proof.
\end{proof}

\end{document}
